%%==================================================
%% chapter04.tex for Doctoral Thesis
%% Based on user input - Images and Tables Commented Out
%%==================================================

% !TeX root = ../main.tex

\chapter{基于隐式神经场的高保真化学资产重建}
\enchapter[High-Fidelity 3D Reconstruction via Implicit Neural Fields]{High-Fidelity 3D Reconstruction of Chemical Apparatus Based on Implicit Neural Fields}
\label{chap:high_fidelity_reconstruction}

本章围绕化学资产重建这一目标，针对透明物体、小物体难以高质量重建和隐式重建结果难以直接用于下游仿真平台的问题提出了一种基于隐式神经场 NeuS 的高保真化学资产重建算法。该算法利用多视角 RGB 图像恢复物体几何，并构建了一个包含数据预处理、隐式场训练、网格提取与优化和纹理贴图生成的完整重建流程。
在透明物体场景下，传统 NeuS 直接以零等值面作为目标表面的做法容易失效，因为透明区域的最大渲染权重并不总与零交叉位置一致，从而导致表面缺失。针对这一问题，本章从透明与不透明区域混合存在的几何特性出发，采用基于绝对距离场局部极小值搜索的统一表面提取策略，实现了对表面的稳定重建。
此外，考虑到小物体在输入图像中像素占比较低、监督信号稀疏，NeuS 训练过程容易长时间停留在过渡态，本章进一步提出学习态增强策略，以加速训练初期的形状收敛，并为后续高频几何细节的刻画预留更多有效迭代。
在纹理重建方面，针对 NeuS 隐式表示难以直接满足仿真平台对显式纹理资产需求的问题，本章构建了一条从 NeuS 隐式颜色场到高分辨率纹理贴图的自动化生成算法，通过 Xatlas 完成 UV 参数化，结合光栅化烘焙、形态学修复与 KNN 补全生成高精度纹理贴图。
实验结果表明，本章算法能够有效提升小物体和透明物体的几何重建质量，并且在纹理贴图生成方面优于 MeshLab 工业通用方案，从而为后续化学实验仿真、视觉感知与机器人操作提供高保真的三维资产基础。

\section{引言}
\ensection{Introduction}


在面向化学实验的具身智能任务场景中，机器人不仅需要识别实验物体的类别，还需要获得其准确的三维几何与纹理信息，从而支撑后续的抓取、操作、仿真与场景交互。化学实验物体本身具有显著的类别、材质和尺度多样性。按照光学属性划分，实验物体既包括烧杯、试管、量筒、冷凝管等透明玻璃器皿，也包括铁架台、滴定管夹、药匙、橡胶塞、试管架等不透明器材。进一步来看，透明物体又可以分为闭合容器和薄壳结构，不透明物体则可以根据形态与使用方式进一步分为常规尺度物体和小尺度物体（小物体）。上述分类体系如图~\ref{fig:chemical_object_taxonomy} 所示。由此可见，面向化学实验的资产重建方法不能只针对单一材质或单一尺度的对象，而需要在统一框架下同时处理透明与不透明物体，并尽可能恢复可用于下游仿真和交互的高质量几何与外观信息。

\begin{figure}[!t]
  \centering
  \includegraphics[width=0.9\textwidth]{figures/ch4/chemical_object_taxonomy}
  \caption{化学实验物体分类体系}
  \label{fig:chemical_object_taxonomy}
\end{figure}

近年来，以 NeRF 为代表的神经隐式重建方法为高保真三维建模提供了新的技术路径。与传统点云、体素和三角网格等显式表示不同，神经隐式场通过连续函数刻画空间中的几何和外观信息，并借助体渲染建立多视角图像与三维空间之间的优化关系，能够在多视角图像监督下恢复较为细致的三维结构和视角相关颜色。对于化学实验中常见的透明玻璃器具，传统多视角几何方法和深度传感方案容易受到透射、折射和反射等光学现象影响，从而产生深度缺失、边界失真和结构不稳定等问题。Ichnowski 等人的工作\cite{ichnowski2022dexnerf}表明，NeRF 由于具备视点相关外观建模能力，能够在透明物体场景中学习与透明性相关的几何信息，并恢复出可用于机器人抓取的深度结果。这说明神经隐式重建方法在透明物体建模方面相较传统方法具有一定优势，也为透明与不透明物体的统一重建提供了可行思路。其中，NeuS\cite{wang2021neus} 将符号距离场与体渲染相结合，能够在多视角 RGB 图像输入下恢复连续、平滑且具有较好几何精度的物体表面。由于 NeuS 同时具备隐式几何表示和体渲染优化能力，可以在统一框架下处理不透明物体的表面重建与透明物体的外观建模，因此适合作为本章化学实验资产重建的基础框架。

但直接将 NeuS 算法应用于化学实验资产仍然面临若干问题。首先，化学实验物体中存在大量小物体，这类物体在图像中占据区域较小，随机采样得到的训练光线多数落在背景区域，导致训练初期几何收敛速度较慢。其次，NeuS 的输出本质上仍然是隐式场，而机器人仿真平台和渲染引擎通常需要显式三角网格、标准 UV 坐标和高分辨率纹理贴图，因而需要进一步完成从隐式表示到显式纹理资产的转换。最后，透明玻璃器具作为化学实验场景中的重要物体类型，其表面成像机理更为复杂，基于零等值面的表面提取方式在该类物体上可能存在表面缺失或结构不稳定等问题。




针对上述问题，本章围绕化学实验资产的高保真重建需求，构建了一套基于 NeuS 的完整三维资产重建流程。整体流程如图~\ref{fig:neus_pipeline} 所示，该流程以多视角 RGB 图像作为输入，依次完成数据预处理、隐式场训练、显式网格提取与纹理贴图生成，最终输出可用于仿真平台和渲染引擎的显式三维资产。

\begin{figure}[htbp]
  \centering
  \includegraphics[width=\textwidth]{figures/ch4/Neus_Pipeline}
  \caption{基于 NeuS 的高保真化学资产重建流程图。包含数据预处理、网络训练、网格提取和纹理贴图生成四个模块。}
  \label{fig:neus_pipeline}
\end{figure}

具体而言，本章首先设计面向 NeuS 重建的数据预处理流程，通过 COLMAP\cite{schonberger2016sfm} 恢复多视角图像的相机参数，并结合主体掩膜完成稀疏点云去噪与空间归一化，为后续隐式场训练建立稳定的坐标参考。在网络训练阶段，本章以 NeuS 作为基础重建框架，并针对化学实验小物体像素占比较低、有效监督信号不足的问题，提出学习态增强策略，以加速训练初期的粗几何收敛，并为后续表面细节学习预留更多有效迭代。在显式资产生成阶段，本章进一步构建从 NeuS 隐式颜色场到高分辨率纹理贴图的自动化生成流程，通过网格提取、UV 参数化、光栅化烘焙、形态学修复和 KNN 补全，生成可直接用于下游仿真平台的纹理资产。此外，考虑到透明玻璃器具在化学实验资产中的重要性，本章在后续实验部分对透明物体的重建结果和局限性进行分析，讨论 NeuS 框架在特殊材质物体重建中的适用性与进一步改进空间。

\section{算法设计与实现}
\ensection{Algorithm Design and Implementation}

\subsection{数据预处理}
\ensubsection{Data Preprocessing}

在基于神经隐式场的三维重建算法中，网络训练主要依赖多视角图像、主体掩膜（Mask）以及相机参数，以建立图像像素与三维空间采样点之间的对应关系。对于 NeuS 算法，还需要对场景空间范围进行合理归一化，以保证射线采样与网络优化过程的数值稳定性。因此，在进行网络训练之前，需要利用多视图几何算法对输入图像进行预处理，从而恢复相机位姿，并为后续隐式场重建建立稳定的坐标参考。

基于上述需求，本节介绍用于高保真资产重建的数据预处理流程。该流程通过多视角图像采集、主体分割以及结构恢复等步骤，为后续神经隐式场网络训练提供准确的相机参数与包围盒。

首先，针对目标物体拍摄多视角的RGB图像序列，采用SFM（Structure-from-Motion\cite{schonberger2016sfm}, SfM）估计相机位姿。在具体实现中，本章采用 COLMAP 开源软件。COLMAP 首先检测图像中的 SIFT 特征点，并在不同视角之间建立特征对应关系；随后利用对极几何约束剔除误匹配点，保证跨视图匹配结果的一致性；在此基础上进一步恢复每张图像的内参与外参，并生成场景的稀疏点云。

% \begin{figure}[htbp]
%     \centering
%     \includegraphics[width=\textwidth,height=3cm]{figures/ch4/colmap_pipeline.png}
%     \caption{COLMAP 多视图重建流程图。该流程主要包括特征提取、特征匹配、几何验证以及稀疏三维重建等步骤。}
%     \label{fig:colmap_pipeline}
% \end{figure}

由于 COLMAP 生成的原始稀疏点云中夹杂较多偏离目标物体区域的离群噪声点，若直接用此进行后续训练，则会导致物体中心、包围范围及尺度估计出现偏差，进而影响射线采样范围的设置以及网络训练的稳定性。如图~\ref{fig:neus_noise} 所示，COLMAP 初始生成的点云往往夹杂较多环境噪声，因此在启动 NeuS 网络训练前，需要对稀疏点云进行精细化的噪声剔除与空间范围修正。

因此本章提出一种基于主体Mask的去噪算法，实现稀疏点云的自动去噪。具体而言，利用物体分割算法\cite{kirillov2023segany}对多视图提取主体 Mask，并将三维空间点根据估计得到的相机参数投影至所有可见视角，统计该点在对应Mask内被判定为主体的频率。若某点在超过半数的可见视图中被判定为背景，则将其识别为噪声点并予以剔除。该算法能够有效清除目标物体周围的散乱噪声，保留能够真实反映目标物体空间范围的有效稀疏点云，为后续场景归一化提供可靠依据。

在完成点云清洗后，为满足 NeuS 算法对输入空间尺度稳定性的要求，本章对剩余点云执行归一化处理。具体而言，首先根据清洗后的稀疏点云估计目标物体的空间中心，并将其平移至世界坐标系原点附近；随后根据点云的整体空间范围计算统一的缩放系数，将场景映射到单位球邻域内，从而为隐式场重建建立一个尺度一致、数值稳定的标准坐标空间，并提升后续体渲染优化与表面重建的稳定性。

\begin{figure}[htbp]
    \centering
    \includegraphics[width=0.8\textwidth]{figures/ch4/Neus_noise.png}
    \caption{基于主体Mask的去噪框架图。COLMAP 生成的原始稀疏点云，包含大量环境噪声。利用重投影到主体Mask中剔除噪声，恢复出干净的点云。}
    \label{fig:neus_noise}
\end{figure}


\subsection{网络训练}
\ensubsection{Network Training}
\label{sec:ch4_network_training}

经过数据预处理后，输入多视角图像已经获得对应的相机内参、相机外参、主体 Mask 以及归一化后的场景空间范围。在此基础上，本节进一步介绍基于 NeuS 的网络训练过程。网络训练是整个重建流程中的核心阶段，其目标是在多视角 RGB 图像监督下学习目标物体的隐式几何场与颜色场，为后续网格提取和纹理贴图生成提供连续的三维表示。

本章以 NeuS 作为基础训练框架，从优化目标来看，网络训练主要由颜色重建损失、主体 Mask 约束损失和距离场正则项共同约束。颜色重建损失用于约束网络预测结果与输入图像之间的一致性，使颜色场能够学习目标物体在不同视角下的外观变化。主体 Mask 约束损失用于限制有效重建区域，使网络更加关注目标物体所在空间，减少背景区域对隐式几何学习的干扰。距离场正则项用于约束符号距离场满足局部连续性，使 SDF 梯度能够稳定反映表面附近的几何变化方向。上述训练目标可以概括为
\begin{equation}
    \text{Loss} =
    \text{Loss}_{\mathrm{color}}
    + \lambda_{\mathrm{mask}}\text{Loss}_{\mathrm{mask}}
    + \lambda_{\mathrm{eik}}\text{Loss}_{\mathrm{eik}},
\end{equation}
其中，$\text{Loss}_{\mathrm{color}}$ 表示颜色重建损失，$\text{Loss}_{\mathrm{mask}}$ 表示主体 Mask 约束损失，$\text{Loss}_{\mathrm{eik}}$ 表示距离场正则项，$\lambda_{\mathrm{mask}}$ 和 $\lambda_{\mathrm{eik}}$ 分别表示对应损失项的权重系数。

上述损失共同决定了隐式几何场的收敛过程。其中，颜色重建损失和主体 Mask 约束损失为目标表面位置提供主要监督，距离场正则项则保证几何场在优化过程中保持较好的连续性和稳定性。对于普通尺度物体，输入图像中通常包含较充分的主体区域，网络能够从颜色一致性和主体区域约束中获得较密集的几何监督，从而较快完成从初始化状态到粗略几何的过渡。然而，当目标物体在图像中占据区域较小时，来自主体区域的有效监督信号会明显减少。此时，距离场正则项虽然能够维持 SDF 的局部平滑性，但难以单独提供足够明确的表面位置引导，从而导致训练初期的几何收敛速度下降。

在化学实验资产重建中，存在不少小物体，它们在图像中仅占据较少像素，例如药匙、橡胶塞、夹具和小型连接件等。对于这类目标，随机采样得到的训练光线中，大量光线落在背景区域或自由空间区域，真正穿过物体表面及轮廓附近的有效光线比例较低，导致训练初期的几何收敛速度较慢。如图~\ref{fig:small_object_pixels} 所示，小物体的有效监督区域较小，网络更容易长时间停留在粗几何尚未稳定的阶段。

\begin{figure}[!t]
\centering
\begin{subfigure}[b]{0.32\textwidth}
\centering
\includegraphics[width=\textwidth]{figures/ch4/small_object_1}
\caption{}
\end{subfigure}
\hfill
\begin{subfigure}[b]{0.32\textwidth}
\centering
\includegraphics[width=\textwidth]{figures/ch4/small_object_2}
\caption{}
\end{subfigure}
\hfill
\begin{subfigure}[b]{0.32\textwidth}
\centering
\includegraphics[width=\textwidth]{figures/ch4/small_object_3}
\caption{}
\end{subfigure}
\caption{小物体示意图。小物体仅覆盖图像中的较小区域，导致有效监督信号较为稀疏。}
\label{fig:small_object_pixels}
\end{figure}

从训练过程来看，NeuS 的几何演化通常可以理解为三个阶段。如图~\ref{fig:neus_three_stage} 所示，网络首先从初始化的单位球出发，随后在过渡态中逐步收缩为目标物体的粗略几何，最后进入学习态以刻画表面细节。对于小物体而言，由于训练早期能够提供有效几何信息的光线数量有限，网络往往需要更多迭代才能完成从初始单位球到粗略物体形状的收敛。若大量训练预算被消耗在过渡态中，后续用于细节学习的有效迭代就会减少，从而影响最终几何质量。

\begin{figure}[htbp]
\centering
\includegraphics[width=\textwidth]{figures/ch4/NeuS_three_tai.png}
\caption{NeuS 训练过程中的几何演化图。训练通常遵循从初始单位球到粗略几何，再到精细表面刻画的演化路径。}
\label{fig:neus_three_stage}
\end{figure}

针对上述问题，本章在网络训练阶段引入学习态增强策略。该策略旨在针对小物体训练早期有效监督不足的问题，通过动态调节采样方向与有符号距离函数梯度方向的权重，引导网络更快完成初始形状收敛，缩短过渡态持续时间，从而尽早进入面向细节优化的学习态，并为高保真几何重建预留更充分的有效迭代空间。

设当前采样点为 $\mathbf{x}$，对应射线方向为 $\mathbf{d}$，有符号距离函数（Signed Distance Function, SDF）记为 $f(\mathbf{x})$。为了刻画当前射线传播方向与局部几何法向之间的关系，首先定义射线方向与 SDF 梯度之间的方向一致性量 $\cos_t$：
\begin{equation}
\cos_t = \mathbf{d} \cdot \nabla f(\mathbf{x}),
\end{equation}
其中，$\mathbf{d}$ 表示当前光线的方向向量，$\nabla f(\mathbf{x})$ 表示该点处的 SDF 梯度。从几何意义上看，$\nabla f(\mathbf{x})$ 对应隐式表面在该点附近的局部法向方向，因此 $\cos_t$ 反映了当前采样方向与局部几何走势之间的一致程度。

然而，在小物体重建过程中，训练早期的几何形状尚不稳定，若过早施加过强的方向约束，容易影响粗几何的快速收敛；而在训练后期，为了提升表面定位精度与细节恢复能力，又需要更严格的方向一致性约束。基于这一考虑，本章进一步引入模式切换因子 $\cos_a$，用于根据训练进度在宽松引导与严格约束之间进行平滑过渡：
\begin{equation}
\cos_{a}=
\begin{cases}
1, & \text{if } \text{cos\_step\_end}=0,\\[3pt]
\min\!\left(1,\dfrac{\text{cur\_step}}{\text{cos\_step\_end}}\right), & \text{otherwise},
\end{cases}
\end{equation}
其中，$\text{cur\_step}$ 为当前训练步数，$\text{cos\_step\_end}$ 为控制阶段切换的超参数。当训练处于初期时，$\cos_a$ 较小；随着训练推进，$\cos_a$ 逐步增大并最终趋近于 $1$，表示优化重点由粗几何收敛逐渐转向细节刻画。

在此基础上，将训练阶段信息与方向约束共同结合，定义自适应采样因子 $\cos_i$：
\begin{equation}
\cos_i=(1-\cos_a)\cos_i^{\text{soft}}+\cos_a \cos_i^{\text{hard}},
\end{equation}
其中，$\cos_i \leq 0$，$\cos_i^{\text{soft}}$ 和 $\cos_i^{\text{hard}}$ 分别表示宽松模式与严格模式下的方向项。由此，$\cos_i$ 在训练初期更接近 $\cos_i^{\text{soft}}$，此时方向约束相对平缓，有利于利用 SDF 梯度先验引导隐式表面快速收缩；而在训练后期，$\cos_i$ 逐渐过渡到 $\cos_i^{\text{hard}}$，从而对采样方向施加更严格的几何一致性约束，以提升表面定位精度。

最后，将上述自适应采样因子用于当前采样区间两端的 SDF 估计，得到
\begin{equation}
\begin{cases}
    \text{SDF}_{\text{next}} = \text{SDF} + \cos_i \cdot D,\\
    \text{SDF}_{\text{prev}} = \text{SDF} - \cos_i \cdot D,
\end{cases}
\end{equation}
其中，$D$ 表示当前采样区间长度。该估计方式的核心思想在于：利用 $\cos_i$ 对相邻位置的 SDF 变化幅度进行自适应调节。当当前区域的几何趋势较明确时，$\cos_i$ 能够放大区间两端的 SDF 差异，从而加快表面位置的收敛；而在几何不确定或表面尚未稳定的区域，较为保守的 $\cos_i$ 则有助于抑制过大的更新步幅，避免跨越真实表面。通过这种方式，本章算法在训练初期能够更快完成由初始单位球到目标粗几何的过渡，而在训练中后期则能够将更多优化预算集中于小物体表面细节的恢复与精细建模。



\subsection{网格提取}
\ensubsection{Mesh Extraction}
\label{sec:ch4_mesh_extraction}

在完成 NeuS 隐式场训练后，需要将网络中的连续几何表示转换为显式三角网格，从而为后续纹理贴图生成、仿真平台导入和机器人交互提供可直接使用的三维资产。

由于化学实验物体在材质属性上存在明显差异，不透明物体与透明物体在隐式场中的表面位置并不完全一致，因此本文在网格提取阶段根据物体材质类型采用不同的表面提取策略。对于不透明物体，采用基于零等值面的标准网格提取流程。对于透明物体，则采用面向透明表面的局部极小值搜索方法，以缓解零等值面提取带来的表面缺失问题。

对于不透明物体，NeuS 学得的符号距离场通常能够较稳定地刻画物体表面。在这种情况下，目标表面可以近似表示为距离场的零等值面，因此本文采用 Marching Cubes\cite{lorensen1987marchingcubes} 算法从训练完成的 SDF 场中提取初始三角网格,并进一步采用 Loop Subdivision\cite{loop1987smooth} 算法进行拓扑细化。该算法通过在每条边上插入新顶点并重计算原始顶点位置，将每个三角形面片分裂为四个子三角形。这一操作能够有效缓解局部面片分布不均的问题，为后续基于顶点的颜色采样提供更为均匀且密集的离散点集。

在获得细化后的高精度网格后，本章进一步利用 NeuS 学习到的颜色场 $c_\theta(\mathbf{x}, \mathbf{d})$ 为几何模型赋予颜色信息。由于原始颜色场具有一定的视角相关性，而工业标准的纹理映射通常要求表面属性尽可能保持视角无关，本章采用一种基于法线方向查询的采样策略：对于每个网格顶点 $\mathbf{v}_i$，根据其邻接面片的法向量计算单位法向 $\mathbf{n}_i$，并将该方向作为查询方向 $\mathbf{d}_i$ 传入颜色场，即
\[
\mathbf{c}_i = c_\theta(\mathbf{v}_i, \mathbf{d}_i), \qquad \mathbf{d}_i = \mathbf{n}_i .
\]
该策略能够在较大程度上保留 NeuS 所学习到的局部纹理细节，并使顶点颜色更符合后续显式纹理映射的需要。

对于透明物体，由于透明物体的成像过程受到透射、折射和反射等因素影响，NeuS 体渲染权重的最大值不一定与 SDF 的零交叉位置严格对齐，直接使用零等值面进行网格提取并不总是可靠。Zhang等人的工作\cite{zhang2024alphaneus} 表明，在透明区域，透明表面更可能对应距离场中的非负局部极小值。如果仍然将零等值面作为唯一的表面提取条件，则容易在透明区域产生表面缺失、局部空洞或结构不连续等问题。



\begin{algorithm}[!t]
\caption{Surface Extraction for Transparent Objects}
\label{algo:transparent_surface_extraction}
\begin{algorithmic}[1]
\STATE \textbf{Stage I: Initial Enclosing Mesh Extraction}
\STATE $f^{a}(\mathbf{x}) \gets |f(\mathbf{x})|$
\STATE $(V,F) \gets \mathrm{MC}(f(\mathbf{x}), r)$

\STATE
\STATE \textbf{Stage II: Coarse Projection}
\FOR{$t=1$ to $T_1$}
    \STATE $C \gets \mathrm{Centroid}(V,F)$
    \STATE Compute the coarse loss
    \[
    L_{\mathrm{coarse}}
    \gets
    \sum_{v_i \in V} f^{a}(v_i)
    +
    \sum_{c_i \in C} f^{a}(c_i)
    +
    \lambda_1 \sum_{v_i \in V} w(v_i)\left\|\Delta v_i\right\|^2
    \]
    \STATE Update $V$ by one optimization step on $L_{\mathrm{coarse}}$
\ENDFOR
\STATE $C^{(1)} \gets \mathrm{Centroid}(V,F)$
\STATE $N^{(1)} \gets \mathrm{FaceNormal}(V,F)$

\STATE
\STATE \textbf{Stage III: Refinement}
\FOR{$t=1$ to $T_2$}
    \STATE $C \gets \mathrm{Centroid}(V,F)$
    \STATE Compute the refinement loss
    \[
    L_{\mathrm{fine}}
    \gets
    \sum_{v_i \in V} f^{a}(v_i)
    +
    \sum_{c_i \in C} f^{a}(c_i)
    +
    \lambda_2 \sum_{c_i \in C}
    \left\|
    \bigl(c_i-c_i^{(1)}\bigr)\times n_i^{(1)}
    \right\|
    \]
    \STATE Update $V$ by one optimization step on $L_{\mathrm{fine}}$
\ENDFOR

\STATE $M^\star \gets (V,F)$

\end{algorithmic}
\end{algorithm}



因此，针对透明物体，本文不再将零等值面简单视为唯一目标表面，而是将与最大渲染权重对齐的位置视为无偏表面，采用了如算法~\ref{algo:transparent_surface_extraction} 所示的流程进行提取。具体实现上，首先对 NeuS 学得的距离场取绝对值，构造绝对距离场，使不透明区域的零等值面和透明区域的非负局部极小值都转化为绝对距离场中的局部极小值。随后采用“先包围、后投影”的策略，先以非零阈值提取包围网格，再通过粗投影和法向约束细化将网格逐步收缩到绝对距离场的局部极小值位置\cite{Hou2023DCUDF}，实现对透明表面的更完整且结构稳定的提取。

\subsection{纹理贴图生成}
\ensubsection{High-Resolution Texture Asset Generation Algorithm}

\begin{figure}[!t]
\centering
\includegraphics[width=\textwidth]{figures/ch4/Xatlas_Stage_PK.png}
\caption{纹理贴图生成流程图。主要包含网格提取、网格颜色查询、UV 坐标生成、光栅化烘焙以及纹理后处理修复五个阶段。}
\label{fig:xatlas_stage_pk}
\end{figure}

在机器人仿真平台~\cite{kolve2017ai2thor,savva2019habitat,shen2021igibson,gu2023maniskill2,li2024chemistry3d,li2025labutopia}中，场景中的物体通常需要以显式三维资产的形式存在，即三维资产不仅要包含显式几何网格，还需要具备材质与纹理信息，以支持渲染、碰撞检测等功能。然而，NeuS 作为一种隐式神经表面重建算法，尽管能够生成高质量的几何表面与颜色信息，但其几何与颜色主要编码在神经网络参数中，难以直接应用到主流机器人仿真平台中。


针对这一问题，本章提出了针对 NeuS 神经隐式场到显式纹理资产的自动化生成流程，其核心思路如图~\ref{fig:xatlas_stage_pk}和流程\ref{algo:neus_texture_generation_en} 所示。

\begin{algorithm}[htbp]
\caption{High-Resolution Texture Asset Generation Based on NeuS}
\label{algo:neus_texture_generation_en}
\begin{algorithmic}[1]

\STATE $M \gets \text{LoopSubdivide}(M)$

\FOR{each vertex $\mathbf{v}_i \in M$}
    \STATE $\mathbf{d}_i \gets \text{EstimateVertexNormal}(\mathbf{v}_i)$
    \STATE $\mathbf{c}_i \gets c_{\theta}(\mathbf{v}_i,\mathbf{d}_i)$
\ENDFOR

\STATE $\mathrm{uv} \gets \text{GenerateUV}(M)$
\STATE $I_{\text{tex}} \gets \text{InitTextureImage}()$

\FOR{each face $f_j \in M$}
    \STATE $\Omega_j \gets \text{FindFaceToUV}(f_j,\mathrm{uv})$
    \FOR{each pixel $p \in \Omega_j$}
        \STATE $I_{\text{tex}}(p) \gets \text{BarycentricInterpolate}(p,f_j,\mathrm{uv})$
    \ENDFOR
\ENDFOR

\STATE $I_{\text{tex}} \gets \text{DilateBoundaryRegions}(I_{\text{tex}})$

\FOR{each uncolored pixel $p \in I_{\text{tex}}$}
    \STATE $I_{\text{tex}}(p) \gets \text{KNNColorInterpolate}(p,I_{\text{tex}})$
\ENDFOR

\end{algorithmic}
\end{algorithm}

在第~\ref{sec:ch4_mesh_extraction} 节中，完成顶点颜色提取后，需要进一步建立三维表面与二维纹理空间之间的映射关系，即生成纹理坐标（UV）。为此，本章引入 Xatlas~\cite{young_xatlas}开源纹理生成库对显式网格进行自动参数化处理，其核心过程如图~\ref{fig:xatlas_uv} 所示。该算法首先依据网格的曲率分布与局部几何特征，对复杂表面进行自动化分块（Charts），再通过基于几何代价的切割与布局优化，将这些局部表面展开并紧凑排列到统一的二维纹理坐标空间中。该过程为后续将三维表面的颜色信息映射到二维纹理图像奠定了几何基础。

\begin{figure}[htbp]
\centering
\includegraphics[width=0.8\textwidth]{figures/ch4/Xatlas_UV.png}
\caption{UV 坐标生成核心流程。}
\label{fig:xatlas_uv}
\end{figure}

在确定 UV 映射关系后，重建流程进入颜色烘焙与光栅化阶段。本章在纹理空间中对每个三角面片执行光栅化操作，并利用重心坐标插值将面片三个顶点上的颜色连续映射到对应的二维像素区域。通过这一过程，原本定义在离散网格顶点上的颜色信息被转化为连续分布在纹理平面中的像素颜色，从而实现了从顶点着色到标准二维纹理贴图的转换。然而，由于光栅化精度限制以及 UV 分块边界处的对齐误差，初步生成的纹理贴图中往往仍会残留少量未被赋色的像素区域，表现为局部的微小黑洞或纹理缺口。

\begin{figure}[htbp]
\centering
\includegraphics[width=0.8\textwidth]{figures/ch4/Xatlas_PNG.png}
\caption{纹理贴图生成核心流程。}
\label{fig:xatlas_png}
\end{figure}

针对上述光栅化后残留的局部空洞问题，本章进一步设计了纹理修复步骤。如图~\ref{fig:xatlas_png} 所示，首先通过形态学膨胀操作引导边缘颜色向邻域适度扩展，以减轻小尺度边界裂缝带来的纹理断裂；随后，对仍未被赋色的像素应用 KNN 搜索，在纹理平面中检索距离该空洞像素最近的有效采样点，并基于邻域颜色对其进行插值补全。通过上述处理，本章能够有效修复 UV 分块边界附近和局部采样不足区域所产生的纹理缺失，从而最终生成连续性更好、可直接应用于仿真平台与图形学渲染系统的高分辨率纹理贴图。


\section{实验结果与分析}
\ensection{Experimental Results and Analysis}

\subsection{实验设置}
\ensubsection{Experimental Setup}

本章实验在 Ubuntu 20.04 操作系统下完成，采用 NVIDIA GeForce RTX 3090 显卡进行训练与推理，CUDA 版本为 11.7，Python 版本为 3.7。

在数据集选择方面，目前公开数据集中缺少面向化学实验资产重建的数据集。理想的化学资产重建数据集需要覆盖常规实验器材、实验小物体与透明玻璃器具，并提供多视角 RGB 图像、相机参数、主体掩膜以及真值几何模型。然而，现有广泛应用于 NeRF 三维重建任务的公开数据集\cite{mildenhall2020nerf,jensen2014large,knapitsch2017tanks,yao2020blendedmvs}侧重于日常物体，难以完整覆盖上述需求。因此，本文根据各实验目标分别选取具有符合对应问题的数据集进行验证。

在对比设置方面，由于本章的所有工作均建立在原始 NeuS 框架之上，因此本章统一采用原始 NeuS 作为对比基线算法，验证各改进模块的有效性。对于小物体实验，基线算法采用原始 NeuS 的训练机制，本章算法在此基础上引入学习态增强策略，以比较两种算法在相同训练预算下的收敛速度、轮廓成形过程以及最终重建质量。对于纹理重建实验，本章将同一组顶点带有颜色的显式网格分别输入 MeshLab \cite{cignoni2008meshlab} 与本章算法进行纹理生成，MeshLab 采用其内置的Texture from Vertex Color功能完成纹理展开与贴图生成。本章算法则采用 Xatlas 进行 UV 参数化，并结合光栅化烘焙、纹理贴图修复，以比较不同纹理生成方案在表面连续性与重渲染保真度上的表现。对于透明物体实验，基线算法直接采用零等值面提取方式生成显式网格，本章算法则采用基于绝对距离场局部极小值搜索的统一表面提取策略，以比较两者在透明区域表面恢复能力上的差异。

在训练与重建参数上，除特殊说明外，各组实验均采用一致的配置。NeuS 的总训练迭代次数设置为 20 万次，Marching Cubes 提取网格时的体素分辨率设置为 $512 \times 512 \times 512$，最终生成的纹理贴图分辨率设置为 $2048 \times 2048$。

\subsection{几何重建实验分析}
\ensubsection{Reconstruction Results and Convergence Analysis on Small Objects}


\begin{figure}[htbp]
    \centering
    \includegraphics[width=\textwidth]{figures/ch4/Neus_Small_Object_PK.png}
    \caption{小物体重建阶段结果可视化对比。}
    \label{fig:neus_small_object_pk}
\end{figure}

本节首先通过定性对比验证算法对小物体的重建能力。如图~\ref{fig:neus_small_object_pk} 所示，在同等计算预算下，本章算法得到的几何重建结果显著优于原始 NeuS。小物体由于在图像空间占据的像素比例较低，在训练初期往往面临监督信号稀疏的问题，导致模型长时间停留在低频特征的学习上。针对这一挑战，本章引入了采样方向与 SDF 梯度方向的动态调节机制，加速网络从初始态向过渡态（10k~50k 次迭代）的跨越。这种策略使得网络能够快速建立粗几何轮廓，从而在后续的学习态（50k~200k 次迭代）中，将更多的计算资源集中于高频细节的精细刻画。

\begin{figure}[!t]
    \centering
    \begin{subfigure}[b]{0.32\textwidth}
        \centering
        \includegraphics[width=\textwidth]{figures/ch4/psnr_curve.png}
        \caption{PSNR}
        \label{fig:psnr_curve_sub}
    \end{subfigure}
    \hfill
    \begin{subfigure}[b]{0.32\textwidth}
        \centering
        \includegraphics[width=\textwidth]{figures/ch4/ssim_curve.png}
        \caption{SSIM}
        \label{fig:ssim_curve_sub}
    \end{subfigure}
    \hfill
    \begin{subfigure}[b]{0.32\textwidth}
        \centering
        \includegraphics[width=\textwidth]{figures/ch4/lpips_curve.png}
        \caption{LPIPS}
        \label{fig:lpips_curve_sub}
    \end{subfigure}
    \caption{小物体重建过程中PSNR、SSIM 和 LPIPS 随训练迭代步数的变化曲线对比。}
    \label{fig:small_object_convergence_curves}
\end{figure}


为了客观评估模型的收敛效率，图~\ref{fig:small_object_convergence_curves} 展示了 PSNR、SSIM 和 LPIPS 三项核心指标随训练迭代步数的变化曲线。实验结果表明，在训练初期约 20k 次迭代时，两种算法的性能较为接近；随着训练逐步进入过渡阶段和细节学习阶段，本章算法在收敛趋势上表现出更明显的优势。具体而言，PSNR 和 SSIM 提升更快，LPIPS 下降更明显，且最终均与基线算法拉开了稳定差距，这表明学习态增强策略有效加速了几何收敛过程，促使模型更早形成稳定的几何与外观表征，并更早进入后续的细节优化阶段，从而在相同训练预算下获得更好的重建质量。


\begin{figure}[!t]
    \centering
    \includegraphics[width=\textwidth]{figures/ch4/transparent_object_vs_neus.png}
    \caption{透明物体几何提取可视化结果对比。第一列为输入 RGB 图像，第二列为 NeuS 的重建结果，第三列为本章算法的重建结果，第四列为对本章算法手动去除外表面后的内部结构可视化结果。图中黑色的片段是自动法线计算过程中部分法向发生翻转所致，并非几何提取错误。}
    \label{fig:neus_transparent_pk}
\end{figure}
图~\ref{fig:neus_transparent_pk}展示了 NeuS 与本章算法的可视化对比结果，原始 NeuS 在透明区域往往难以得到稳定的零交叉表面，导致最终提取出的网格存在明显的空洞和缺失。相比之下，本章算法通过在绝对距离场上搜索局部极小值，并将透明区域与不透明区域统一纳入同一提取框架，有效缓解了上述问题，并保留较好的内部几何信息。需要说明的是，在结果中可见少量黑色几何片段，这主要是由于自动法向计算过程中部分法向方向发生翻转所致，并非几何提取错误所引起。


表~\ref{tab:transparent_cd_compare} 从定量的角度展示了本章算法与原始 NeuS 在 Synthetic Blender 数据集上的 Chamfer Distance 对比结果。可以看出，本章算法在全部测试对象上均取得了更低的 CD 值，平均值由 23.9\,mm 显著降低至 5.4\,mm。这表明，相比于直接基于零等值面提取网格的方式，本章所提出的统一表面提取策略能够更准确地定位透明物体的真实外表面，从而有效提升几何重建精度。

\begin{table}[!t]
\centering
\caption{透明物体重建的 Chamfer Distance$\downarrow$对比（单位：mm）}
\label{tab:transparent_cd_compare}
\small
\setlength{\tabcolsep}{12pt}
\begin{tabular}{ccc}
\toprule
物体类别 & NeuS & Ours \\
\midrule
Bottle     & 7.1  & \textbf{3.7} \\
Case       & 21.2  & \textbf{5.4} \\
Jar        & 43.2 & \textbf{9.2} \\
Jug        & 10.8  & \textbf{4.1} \\
Snowglobe  & 37.2  & \textbf{4.6} \\
\midrule
平均值     & 23.9 & \textbf{5.4} \\
\bottomrule
\end{tabular}
\end{table}

综合实验分析，针对原始 NeuS 在透明区域中容易因零交叉位置与真实表面不一致而产生表面缺失，从而影响最终网格的完整性与几何准确性问题，本章算法通过在绝对距离场中搜索局部极小值位置，并将透明区域与不透明区域纳入统一的表面提取框架，在透明物体重建任务中表现出更好的稳定性与适用性，能够在保持整体结构完整性的同时，更准确地恢复物体的几何边界。

综合最终重建结果与训练过程分析可以看出，本章所提算法能够有效缓解小物体在训练初期因监督信号稀疏而导致的长时间过渡问题，使模型更早进入稳定的几何细化阶段，从而为后续高频细节学习预留更多有效迭代空间，并最终提升小物体的重建精度与细节完整性。


\subsection{纹理贴图生成实验分析}
\ensubsection{3D Asset Reconstruction Results and Analysis}
\begin{figure}[!t]
  \centering
  \includegraphics[width=0.9\textwidth,height=0.75\textheight]{figures/ch4/texture_example}
  \caption{本章算法生成的三维资产示例。从左到右依次为输入图像、重建得到的三维资产、对应几何模型和纹理贴图。}
  \label{fig:texture_example}
\end{figure}
本节旨在验证所提算法将隐式神经场结果转换为高保真显式三维资产的有效性。图~\ref{fig:texture_example} 展示了本章所提的高保真资产重建算法的重建效果，结果表明，本章算法能够将输入图像中的物体重建为较精细的几何模型和纹理贴图。对于苦瓜表面的凹凸细节和牛角包的烘焙纹理特征，算法均能够较准确地刻画。对于玉米的颗粒重复结构，生成的纹理贴图也较好地保留了原始图像中的外观特征，未出现明显的错乱与遗漏。

\FloatBarrier

为了进一步评价重建资产的外观保真度，本章将所提算法与工业开源软件 MeshLab 进行了定性对比，如图~\ref{fig:meshlab_compare}所示。MeshLab 算法在三角形面片之间存在肉眼可见的纹理断层，而本章算法通过稳定的 UV 参数化与后处理修复策略，生成的纹理贴图在表面连续性和细节还原方面表现更为出色。

\begin{figure}[!t]
  \centering
  \includegraphics[width=0.8\textwidth]{figures/ch4/MeshLab_Compare.png}
  \caption{MeshLab 与本章算法生成的纹理贴图对比。}
  \label{fig:meshlab_compare}
\end{figure}

为进一步量化评估重建模型的外观保真度并验证各核心模块的效果，本章以 PSNR 为评价指标，在 PyTorch3D 中对 21 个物体生成的带纹理 OBJ 模型进行了逐视角重渲染对比, 对同一物体全部视角的 PSNR 取平均，得到该物体的平均 PSNR。实验共设置三种方案：基于 MeshLab 顶点颜色烘焙的基准方案（MeshLab）、剔除形态学修复模块的消融版本（Ours w/o Repair）以及本章完整算法（Ours）。结果如表~\ref{tab:texture_psnr_compare} 所示，本章完整算法在绝大多数场景中均取得了最优表现，其平均 PSNR 达到 25.509 dB，相较于 MeshLab 基准和无修复版本分别提升了 2.574 dB 与 1.432 dB。这一定量结果表明，稳定的 UV 参数化为纹理映射奠定了良好的几何基础，而本章设计的后处理修复算法进一步提升了纹理贴图精度。

综上所述，定性观察与定量实验均表明，本章构建的自动化纹理生成算法能够高效地将 NeuS 的隐式颜色场转换为高质量的显式纹理资产，能够为后续的物理仿真、图形渲染与具身智能任务提供高质量的三维资产基础。


\begin{table}[!t]
\centering
\caption{纹理贴图生成算法的 PSNR$\uparrow$ 对比（单位：dB）}
\label{tab:texture_psnr_compare}
\small
\setlength{\tabcolsep}{5pt}
\begin{tabular}{cccc}
\toprule
物体类别 & MeshLab & Ours w/o Repair & Ours \\
\midrule
Apple & 21.310 & 22.127 & \textbf{23.938} \\
Bitter melon & 23.791 & 24.317 & \textbf{25.179} \\
Banana & 23.882 & 24.001 & \textbf{24.030} \\
Bagel & 26.105 & 26.854 & \textbf{28.766} \\
Carrot & 26.288 & 26.335 & \textbf{28.967} \\
Corn & 24.354 & 25.012 & \textbf{25.711} \\
Sausage roll & 22.421 & 24.981 & \textbf{27.121} \\
Fried chicken & 20.755 & 23.711 & \textbf{26.015} \\
Ukulele & 21.952 & 21.968 & \textbf{24.913} \\
Hamburg & 24.462 & 25.565 & \textbf{27.586} \\
Lemon & 28.686 & 29.489 & \textbf{30.191} \\
Luffa & 23.758 & 24.899 & \textbf{27.992} \\
Mango & 20.626 & 20.746 & \textbf{21.844} \\
Walnut & \textbf{21.679} & 20.020 & 20.973 \\
Cepa & 26.490 & 26.797 & \textbf{28.786} \\
Blue-and-white vase & 23.698 & 23.698 & \textbf{23.809} \\
Celadon vase & 22.956 & 24.151 & \textbf{25.234} \\
Pineapple & 13.425 & 16.492 & \textbf{18.517} \\
Strawberry cream pastry & 16.824 & 20.046 & \textbf{21.237} \\
Croissant & 21.718 & 24.068 & \textbf{26.238} \\
Chocolate cake & 26.450 & 27.893 & \textbf{28.637} \\
\midrule
平均值 & 22.935 & 24.077 & \textbf{25.509} \\
\bottomrule
\end{tabular}
\end{table}




\subsection{局限性分析}
\ensubsection{Limitations Analysis}

尽管本章算法在透明物体表面提取、小物体几何重建以及纹理贴图生成等方面取得了较好的实验效果，但仍然存在一些局限性，有待在后续工作中进一步完善。

(1)在小物体重建加速方面，本章提出的学习态增强策略主要是在 NeuS 训练框架内，通过调节训练前期的几何收敛过程来提升有限训练预算下的小物体重建效率。然而，这一策略并未从根本上改变 NeuS 训练速度较慢的问题。相较之下，近年来以 Instant-NGP 为代表的高效神经表示算法\cite{muller2022instantngp,Li_2023_CVPR,Wang_2023_ICCV} ，通过多分辨率哈希编码等机制，显著提升了神经场的训练与查询效率，在保持较强表达能力的同时大幅缩短了优化时间。相比这类高效表示，本章算法仍然建立在较为传统的隐式场优化范式之上，因此在大规模数据处理、快速重建和交互式建模等应用场景中仍存在一定限制。未来可以考虑将本章针对小物体训练动态的改进思想，与 Instant-NGP 一类高效编码框架相结合，在保持几何细节恢复能力的同时进一步降低训练成本，从而实现更快、更稳定的小物体高保真重建。

(2)在纹理贴图生成方面，本章构建的自动化算法能够将 NeuS 学得的隐式几何与颜色信息转换为主流渲染引擎可直接使用的显式资产，但这一过程本身不可避免地伴随着一定的信息损失。其根本原因在于，NeuS 算法采用的是连续隐式表示，颜色通常由空间位置与视角方向共同决定，能够在隐式场中保留较丰富的连续外观信息。而当其被转换为三角网格、固定 UV 坐标以及二维纹理贴图后，原始隐式表示中的连续性、视角相关性以及部分高频细节将被压缩到有限分辨率的显式表示中。这意味着从隐式场到显式纹理模型的转换并不是严格无损的表达映射，而是一种近似映射。因此，即使本章算法已经通过光栅化烘焙、形态学修复与 KNN 补全尽可能提升了纹理完整性，最终结果仍难以完全等价于原始隐式表示，不可避免地存在一定的纹理信息损失。未来可以考虑引入基于材质分解的显式表示、视角相关纹理建模，或采用隐式与显式混合表示，以进一步减小这一转换过程带来的纹理信息损失。

(3)在透明物体重建方面，本章算法虽然能够较好地重建透明物体的几何网格，并在一定程度上缓解 NeuS 零等值面提取所带来的表面缺失问题，但其结果仍然存在一定的不稳定性。首先，由于透明区域对应的几何信号较弱，且后续网格法向通常依赖局部邻域自动估计，因此在部分区域仍可能出现法向方向翻转的问题，进而表现为局部黑面、明暗异常或表面朝向不一致。其次，本章算法仍然建立在多视图图像对隐式场的优化基础之上，当场景中存在较强的折射、高光反射或复杂环境光干扰时，透明物体的真实几何边界与图像观测之间的对应关系会进一步减弱，此时隐式场学习到的表面位置可能出现波动，从而影响最终提取结果的稳定性。因此，对于具有复杂光照条件的透明物体，本章算法的鲁棒性仍有提升空间。最后，本文实验主要验证了在 NeuS 框架下改进模块的有效性，尚未探索超越 NeuS 框架或结合其他先进透明物体重建方法的潜力。未来研究可在此基础上考虑多框架融合、物理光学建模或跨方法集成，以进一步提升透明物体的重建精度和鲁棒性。

\section{本章小结}
\ensection{Summary}

本章面向化学实验场景中透明物体、小物体难以高质量重建和机器人仿真平台对显式资产的需求，提出了一种基于隐式神经表面 NeuS 的高保真化学资产重建算法。围绕高保真化学资产重建这一目标，本章从数据预处理、隐式场训练、网格提取到纹理贴图生成，构建了一条完整的重建流程，为后续化学实验仿真、视觉感知与机器人操作提供可直接使用的三维资产。

在算法上，本章首先设计了面向 NeuS 训练的数据预处理流程。通过多视图图像采集、COLMAP 相机位姿恢复、基于主体Mask的稀疏点云去噪以及空间归一化处理，建立了稳定的相机参数与坐标参考，为后续隐式场优化提供了可靠输入。在NeuS 网络中，本章围绕两个几何关键问题进行了针对性改进：一是在训练阶段，针对小物体训练初期监督信号稀疏、收敛缓慢的问题，提出学习态增强策略，以加速几何初始收敛并为后续细节学习预留更多有效迭代。二是在网格提取阶段，针对透明物体表面与零等值面不稳定对应的问题，采用基于绝对距离场局部极小值搜索的统一表面提取策略，以提高透明区域与不透明区域的几何恢复质量。进一步地，为满足下游仿真平台对显式资产的使用需求，本章构建了从 NeuS 隐式颜色场到高分辨率纹理贴图的自动化生成算法，利用 Xatlas 进行 UV 参数化，并结合光栅化烘焙、形态学修复与 KNN 补全，生成了适用于主流渲染引擎和仿真平台的纹理模型。

实验结果表明，本章算法在几何与纹理两个层面均取得了较好的效果。在透明物体重建中，所提出的表面提取策略能够更完整地恢复器皿外轮廓，并减少内部伪影；在小物体重建中，学习态增强策略能够缩短过渡阶段并提升边缘与细节恢复能力；在纹理生成中，本章算法相比 MeshLab 等通用工具链能够有效减少纹理接缝、空洞与边缘锯齿，获得更加连续和平滑的表面外观。综上，本章提出的算法能够有效提升化学实验资产三维重建的保真度与工程可用性，为后续面向具身智能训练的高质量化学实验场景构建奠定了三维资产基础。

